\subsection{A cardinal bound on the quotients of two-dimensional groups}

\begin{claim}{A.19}{Proved}{Two-dimensional groups with small homology have small detecting quotients}
Let $G$ have a connected CW classifying space $X=BG$ of dimension
at most two. Let $d(H_1(G;\mathbb Z))$ be the least cardinality of
an abelian generating set of $H_1(G;\mathbb Z)$, and put
\[
 \nu=\max\bigl(\aleph_0,d(H_1(G;\mathbb Z)),
                    \operatorname{rank}H_2(G;\mathbb Z)\bigr).
 \tag{A.19.1}
\]
Here $H_2(G;\mathbb Z)$ is free abelian, as follows from the
specified two-dimensional model. For every $g\in G\setminus\{1\}$,
there is a surjective homomorphism
\[
 p_g:G\twoheadrightarrow Q_g,
 \qquad |Q_g|\leq\nu^{\aleph_0},\qquad p_g(g)\ne1.
 \tag{A.19.2}
\]
Thus $G$ is residually in groups of cardinality at most
$\nu^{\aleph_0}$. In particular, countable ordinary integral
homology in degrees one and two gives the bound $2^{\aleph_0}$.
\end{claim}

\begin{proof}
We use the precise coefficient-injectivity assertion of
\cite[Theorem~1]{D_OpenAI2026Nonsingular}: for every group $P$,
every finite family $w_1,\ldots,w_m\in P*F(t_1,\ldots,t_n)$
whose exponent-sum rows are linearly independent over $\mathbb Q$
has an injective canonical coefficient map
\[
 P\longrightarrow
 (P*F(t_1,\ldots,t_n))/\langle\!\langle w_1,\ldots,w_m\rangle\!\rangle.
 \tag{A.19.3}
\]
The source used here is its new, unrefereed October 2026 preprint
version. Its fixed-space incidence argument and its operator and
planar arguments have been directly checked in the accompanying
source audit. No torsion or approximation hypothesis on $P$ is
part of (A.19.3). The rest of the proof constructs all comparison
maps used to obtain (A.19.2).

\paragraph{A small homology core containing the chosen element.}
Fix a based loop representing $g$. As $X$ has no cells above
dimension two, its group $H_2(X;\mathbb Z)$ is a subgroup of the
free abelian cellular group $C_2(X;\mathbb Z)$ and is therefore
free abelian. Choose cellular two-cycles representing a basis of
this homology group, and choose cellular one-cycles representing
an abelian generating set of $H_1(X;\mathbb Z)$. The chosen
cycles, the chosen loop, and connecting paths are contained in
a connected subcomplex $Y_0\subset X$ with at most $\nu$ cells.
This uses only finite supports for each cycle and the fact that
a compact subset of a CW complex is contained in a finite
subcomplex. The maps from $Y_0$ are onto on $H_1$ and $H_2$.

Given a connected subcomplex $Y_n$ with at most $\nu$ cells,
its fundamental group and first homology both have cardinality
at most $\nu$. For each element of
$\ker(\pi_1(Y_n)\to\pi_1(X))$, choose a representing loop and
one nullhomotopy disk in $X$, and add a finite subcomplex
containing its image. For each element of
$\ker(H_1(Y_n;\mathbb Z)\to H_1(X;\mathbb Z))$, choose a
cellular one-cycle and a cellular two-chain in $X$ bounding
it, and add that chain's finite support. Add connecting paths
as needed. The result is a connected $Y_{n+1}$ with at most
$\nu$ cells. Put $Y=\bigcup_{n<\omega}Y_n$.

Every loop and every finite cellular cycle in $Y$ already lies
in one $Y_n$. The specified nullhomotopies consequently show
that the actual inclusion $Y\hookrightarrow X$ induces an
injection on fundamental groups. The specified bounding chains
show that its map on $H_1$ is injective; its initial choice makes
that map surjective. The map on $H_2$ is injective because the
inclusion of cellular two-chains is injective and neither
complex has cellular three-boundaries. Its initial choice also
makes this map surjective. Thus $Y\hookrightarrow X$ is an
ordinary integral homology isomorphism in every degree.

The proved injection on fundamental groups identifies the
universal cover of $Y$ with a component of the preimage of $Y$
in the contractible universal cover of $X$. It is a
two-dimensional subcomplex of that cover. The same cellular
two-cycle injection makes its second homology zero. It is
simply connected and has no homology in higher degrees, so
Hurewicz and Whitehead make it contractible. Hence $Y$ is an
aspherical CW complex with at most $\nu$ cells. Write
$H=\pi_1(Y)$ and $h\in H$ for the chosen element mapping to
$g$. Then $|H|\leq\nu$ and $h\ne1$.

\paragraph{An actual relative cellular isomorphism.}
Choose a maximal tree $T_Y$ of $Y^{(1)}$, and extend it to a
maximal tree $T$ of $X^{(1)}$. Since $T_Y$ already connects
all the vertices of $Y$, the equality $T\cap Y=T_Y$ holds:
an additional edge of $Y$ in $T$ would create a cycle. The
quotient map $q:X\to\overline X=X/T$ restricts to
$Y\to\overline Y=Y/T_Y$, and $\overline Y$ is an actual
subcomplex of $\overline X$. Collapsing a contractible CW
subcomplex is a homotopy equivalence, so both these quotient
maps induce isomorphisms on fundamental groups and homology.
Both quotient complexes have one vertex. We use their
specified isomorphisms to identify the groups with $G$ and $H$.

Let $S$ be the set of oriented one-cells of
$\overline X-\overline Y$, and let $R$ be its set of two-cells.
The relative cellular complex is concentrated in degrees two
and one. Its homology is zero by the preceding homology
isomorphisms. Consequently its actual cellular boundary is
an isomorphism
\[
 \delta:\mathbb Z^{(R)}\xrightarrow{\ \cong\ }
                         \mathbb Z^{(S)}.
 \tag{A.19.4}
\]
Choose for each two-cell $r$ a finite word
$w_r\in H*F(S)$ representing its attaching loop. Choose the
word within the finite subgraph supporting that attaching
map. Let $W(r)$ be the full finite set of relative one-cells
whose interiors meet the attaching image of $r$; in particular it contains
every letter of $w_r$, including letters of exponent sum
zero. Let $I(s)$ be the finite support of $\delta^{-1}(s)$.

Call $(S',R')$, with $S'\subset S$ and $R'\subset R$, closed
when
\[
 r\in R'\ \Longrightarrow\ W(r)\subset S',
 \qquad
 s\in S'\ \Longrightarrow\ I(s)\subset R'.
 \tag{A.19.5}
\]
The first condition selects an actual CW subcomplex together
with $\overline Y$. It also makes $\delta$ carry the selected
two-cell summand into the selected one-cell summand. The
second makes $\delta^{-1}$ preserve these summands in the
reverse direction. Thus their restrictions are inverse
isomorphisms.

Starting with a closed pair, add one unused cell and close
under the two finite-support operations in (A.19.5) for
countably many rounds. At most countably many cells are
added. Increasing unions of closed pairs are closed.
Well-ordering all the cells and repeatedly taking the first
unused cell gives a continuous filtration by actual
subcomplexes
\[
 \overline Y=X_0\subset X_1\subset\cdots\subset
 X_\tau=\overline X
 \tag{A.19.6}
\]
for a small ordinal $\tau$. Each successor pair has at most
countably many cells. If its closed pairs are
$(S_\alpha,R_\alpha)\subset(S_{\alpha+1},R_{\alpha+1})$,
the maps $\delta$ and $\delta^{-1}$ preserve both pairs of
summands. They induce inverse quotient maps
\[
 \delta_\alpha:
 \mathbb Z^{(R_{\alpha+1}\setminus R_\alpha)}
       \xrightarrow{\ \cong\ }
 \mathbb Z^{(S_{\alpha+1}\setminus S_\alpha)}.
 \tag{A.19.7}
\]
The forward map is exactly the relative cellular boundary of
$(X_{\alpha+1},X_\alpha)$.

Put $G_\alpha=\pi_1(X_\alpha)$. Van Kampen gives the precise
successor presentation
\[
 G_{\alpha+1}=
 \bigl(G_\alpha*F(S_{\alpha+1}\setminus S_\alpha)\bigr)
 \big/
 \langle\!\langle \widetilde w_r:
                    r\in R_{\alpha+1}\setminus R_\alpha
                    \rangle\!\rangle,
 \tag{A.19.8}
\]
where $\widetilde w_r$ is the original attaching word with
older letters read as coefficients in $G_\alpha$. Its
exponent-sum matrix in the new variables is the transpose
of (A.19.7). In particular every finite subfamily of its
rows is linearly independent over $\mathbb Q$. At limits,
$G_\beta=\mathop{\rm colim}_{\alpha<\beta}G_\alpha$ by
the CW union form of van Kampen. These presentation and
colimit statements require no unproved injectivity of an
intermediate group map.

\paragraph{A fixed target for all the countable systems.}
We construct a group $E$ of cardinality at most $\nu$,
containing $H$, in which every finite nonsingular system
over $E$ has a solution. Start with $E_0=H$. Given $E_n$,
there are at most $\nu$ finite nonsingular systems with
coefficients in $E_n$: all their words are finite, the
number of variables and equations is finite, and
$|E_n|\leq\nu$. Well-order this set and successively
adjoin the variables and relators of each listed system.
At each step its coefficients map from $E_n$ into the
current group. The exponent rows remain the same under
this coefficient map. Formula (A.19.3) therefore proves
that every resulting coefficient map is injective.
Use group colimits at limit stages. The resulting group
$E_{n+1}$ has cardinality at most $\nu$, since at most
$\nu$ finite sets of generators have been added.

Let $E=\mathop{\rm colim}_{n<\omega}E_n$. All structure
maps into $E$ are injective, and $|E|\leq\nu$. The
finitely many coefficients of any finite nonsingular
system over $E$ belong to one $E_n$. That system was
solved in the construction of $E_{n+1}$, so its solution
maps into $E$.

Now take the explicit reduced power
\[
 U=E^{\mathbb N}/E^{(\mathbb N)},
 \qquad
 E^{(\mathbb N)}=
 \{(e_k):e_k=1\text{ for all but finitely many }k\}.
 \tag{A.19.9}
\]
The subgroup displayed in (A.19.9) is normal: coordinatewise
conjugation preserves its finite support. A nonidentity
constant sequence has infinite support, so the diagonal
map $E\hookrightarrow U$ is injective. In particular its
restriction gives an embedding $j:H\hookrightarrow U$.
Moreover $|U|\leq\nu^{\aleph_0}$.

We check exactly the solution property needed below.
Let $v_1,v_2,\ldots$ be a countable system over $U$ in
countably many variables, every finite subfamily of whose
exponent rows is linearly independent over $\mathbb Q$.
Choose once and for all a representative in $E^{\mathbb N}$
for each coefficient occurring in these words. At coordinate
$k$, read these representatives in $E$ and solve the first
$k$ equations there. This is a finite nonsingular system:
its finite set of occurring variables contains every
nonzero exponent entry of those rows, and changing the
coefficients does not change the exponent matrix. Such a
solution exists by the property of $E$. Give every unused
variable the value $1$ at that coordinate. For each
variable, these coordinate values form an element of
$E^{\mathbb N}$ and hence an element of $U$.
The value of the word $v_i$ in these sequences is $1$ at
every coordinate $k\geq i$. Its value thus belongs to
$E^{(\mathbb N)}$, so $v_i=1$ in $U$. This solves the
entire countable system in this one target. Finite or
empty systems are covered by the same construction,
using all their equations once $k$ is sufficiently large.

\paragraph{The detecting homomorphism.}
We construct homomorphisms $\phi_\alpha:G_\alpha\to U$
compatible with every map in (A.19.6). Set $\phi_0=j$.
At a successor stage apply $\phi_\alpha$ to the
coefficients of the countably many words in (A.19.8).
The exponent rows in the new variables remain unchanged,
so the solution property of (A.19.9) supplies elements
of $U$ for all those variables simultaneously. The
homomorphism from the displayed free product determined
by these elements and $\phi_\alpha$ kills every specified
relator. It therefore descends to
$\phi_{\alpha+1}:G_{\alpha+1}\to U$, with the exact
compatibility
$\phi_{\alpha+1}\circ(G_\alpha\to G_{\alpha+1})
=\phi_\alpha$. At a limit the group-colimit universal
property gives the unique compatible map to $U$.

The terminal map, composed with the specified isomorphism
$q_*:G\to\pi_1(\overline X)$, is a homomorphism
$\Phi_g:G\to U$ whose restriction to the actual subgroup
$H\subset G$ is $j$. Consequently
$\Phi_g(g)=j(h)\ne1$. Let $Q_g=\Phi_g(G)$ and let
$p_g$ be $\Phi_g$ with codomain restricted to this image.
This is the surjection (A.19.2), and
$|Q_g|\leq|U|\leq\nu^{\aleph_0}$ proves its cardinal
bound. If the two homology groups in the statement are
countable, then $\nu=\aleph_0$ and
$\nu^{\aleph_0}=2^{\aleph_0}$, as claimed.
\end{proof}
